\documentclass[11pt, a4paper]{article}

% --- UNIVERSAL PREAMBLE BLOCK ---
\usepackage[a4paper, top=2.5cm, bottom=2.5cm, left=2cm, right=2cm]{geometry}
\usepackage{fontspec}

\usepackage[russian, bidi=basic, provide=*]{babel}

\babelprovide[import, onchar=ids fonts]{russian}
\babelprovide[import, onchar=ids fonts]{english}

% Set default/Latin font to Sans Serif in the main (rm) slot
\babelfont{rm}{Noto Sans}
% Assign a specific font for Russian text
\babelfont[russian]{rm}{Noto Sans}

% Add because main language is not English
\usepackage{enumitem}
\usepackage{booktabs}
\setlist[itemize]{label=-}

\usepackage{listings}
\usepackage{xcolor}

% Настройка внешнего вида листингов кода
\definecolor{codegreen}{rgb}{0,0.6,0}
\definecolor{codegray}{rgb}{0.5,0.5,0.5}
\definecolor{codepurple}{rgb}{0.58,0,0.82}
\definecolor{backcolour}{rgb}{0.95,0.95,0.92}

\lstdefinestyle{mystyle}{
    backgroundcolor=\color{backcolour},   
    commentstyle=\color{codegreen},
    keywordstyle=\color{magenta},
    numberstyle=\tiny\color{codegray},
    stringstyle=\color{codepurple},
    basicstyle=\ttfamily\footnotesize,
    breakatwhitespace=false,         
    breaklines=true,                 
    captionpos=b,                    
    keepspaces=true,                 
    numbers=left,                    
    numbersep=5pt,                  
    showspaces=false,                
    showstringspaces=false,
    showtabs=false,                  
    tabsize=2,
    language=C
}
\lstset{style=mystyle}

\begin{document}

\begin{titlepage}
    \centering
    \vspace*{2cm}
    
    {\large \textbf{МОСКОВСКИЙ АВИАЦИОННЫЙ ИНСТИТУТ}} \\[0.5cm]
    {\large \textbf{(НАЦИОНАЛЬНЫЙ ИССЛЕДОВАТЕЛЬСКИЙ УНИВЕРСИТЕТ)}} \\[2cm]
    
    {\Large \textbf{ОТЧЕТ}} \\
    {\large по лабораторной работе № 2} \\
    {\large по дисциплине «Операционные системы»} \\[1.5cm]
    
    {\Large \textbf{Тема: Управление процессами и межпроцессное взаимодействие (Pipes) в POSIX}} \\[3cm]
    
    \begin{flushright}
        \textbf{Выполнил:} \\
        студент группы М8О-203БВ-25 \\
        Эль-Тахир Роман Райидович \\[1cm]
        
        \textbf{Проверил:} \\
        преподаватель \\
        Миронов Е.С.
    \end{flushright}
    
    \vfill
    {\large Москва -- \the\year}
\end{titlepage}

\newpage

\section{Цель работы}
Изучить механизмы создания процессов и межпроцессного взаимодействия в операционных системах Linux на примере неименованных каналов pipe, перенаправление стандартных потоков ввода-вывода,
а также организацию кода с отделением работы с системными вызовами от общей логики программы.(\texttt{fork}, \texttt{exec}, \texttt{waitpid}) и организации межпроцессного взаимодействия с использованием неименованных каналов (\texttt{pipe}).

\section{Постановка задачи}

Необходимо разработать две \textbf{отдельные} программы: родительскую
(\texttt{parent}) и дочернюю (\texttt{child}).

\begin{itemize}[noitemsep]
  \item Родительский процесс запрашивает у пользователя два имени файлов:
        первая введённая строка~--- файл для \texttt{child1}, вторая~--- для \texttt{child2}.
  \item Родитель создаёт два неименованных канала \texttt{pipe1} и \texttt{pipe2},
        открывает выходные файлы и запускает два дочерних процесса.
  \item Все последующие строки, введённые в консоль, распределяются по
        каналам: \textbf{нечётные} строки (1-я, 3-я, \dots) отправляются в
        \texttt{pipe1} (для \texttt{child1}), \textbf{чётные} (2-я, 4-я, \dots)~---
        в \texttt{pipe2} (для \texttt{child2}).
  \item Дочерние процессы читают строки из своего канала через \texttt{stdin},
        \textbf{инвертируют} каждую строку и выводят в \texttt{stdout}, который
        родитель перенаправил в соответствующий файл.
  \item Программа должна работать в Linux (POSIX).
\end{itemize}

\textbf{Пример.} Ввод: \texttt{out1.txt}, \texttt{out2.txt}, \texttt{abc}, \texttt{def},
\texttt{ghi}, \texttt{jkl}. Результат: в \texttt{out1.txt}~--- \texttt{cba}, \texttt{ihg};
в \texttt{out2.txt}~--- \texttt{fed}, \texttt{lkj}.

\section{Структура проекта и архитектура}

Проект собирается системой CMake и состоит из следующих модулей:

\begin{center}
\begin{tabular}{@{}p{4.2cm}p{10cm}@{}}
\toprule
\textbf{Файл / папка} & \textbf{Назначение}\\
\midrule
\texttt{common/\allowbreak{}errors.c} & Вывод сообщений об ошибках (\texttt{errno}).\\
\texttt{common/\allowbreak{}io\_\allowbreak{}utils.c} & Чтение строки произвольной длины (\texttt{read\_\allowbreak{}line}).\\
\texttt{os/\allowbreak{}os.h} & Единый интерфейс работы с процессами, каналами и файлами.\\
\texttt{os/\allowbreak{}os\_\allowbreak{}linux.c} & Реализация интерфейса для Linux (POSIX).\\
\texttt{parent/\allowbreak{}main.c} & Родительская программа.\\
\texttt{child/\allowbreak{}main.c} & Дочерняя программа (инверсия строк).\\
\texttt{tests/\allowbreak{}run\_\allowbreak{}test.cmake} & Автоматический тест (\texttt{ctest}).\\
\bottomrule
\end{tabular}
\end{center}

Ключевая идея архитектуры~--- \textbf{слой абстракции}. Файл \texttt{os.h}
объявляет набор функций (\texttt{os\_\allowbreak{}pipe\_\allowbreak{}create}, \texttt{os\_\allowbreak{}process\_\allowbreak{}start},
\texttt{os\_\allowbreak{}write\_\allowbreak{}all}, \texttt{os\_\allowbreak{}close}, \texttt{os\_\allowbreak{}process\_\allowbreak{}wait}\dots) и
типы данных. Логика родителя в \texttt{parent/\allowbreak{}main.c}
не содержит прямых вызовов \texttt{fork} или \texttt{execl}: они скрыты
внутри \texttt{os\_\allowbreak{}linux.c}. Типы вынесены в \texttt{os.h}:

\begin{lstlisting}
typedef int os_handle_t;             /* Linux: fd (int)   */
#define OS_INVALID_HANDLE (-1)
\end{lstlisting}

\subsection{Схема взаимодействия процессов}

\begin{verbatim}
                    +--------------------------------+
   клавиатура --->  |            parent              |
   (stdin)          | строка 1,2: имена файлов       |
                    | нечётные --> pipe1             |
                    | чётные   --> pipe2             |
                    +---------+----------+-----------+
                              |          |
                         pipe1|          |pipe2
                              v          v
                    +-----------+     +-----------+
                    |  child1   |     |  child2   |
                    | stdin=p1  |     | stdin=p2  |
                    | stdout=f1 |     | stdout=f2 |
                    +-----+-----+     +-----+-----+
                          v                 v
                      файл 1            файл 2
                   (строки перевёрнуты)
\end{verbatim}
\section{Описание основных системных вызовов (Linux)}

\subsection{Директива \texttt{\#define \_GNU\_SOURCE}}

Макрос, который необходимо указать \textbf{до} всех \texttt{\#include}.
Он включает в заголовочных файлах библиотеки glibc нестандартные
расширения Linux. В нашем коде он нужен для функции \texttt{pipe2()}: без него её
объявления в \texttt{<unistd.h>} не будет.

\subsection{Тип \texttt{os\_handle\_t}}

Собственный тип проекта. Это \texttt{int}~--- файловый дескриптор, то есть
номер открытого файла, канала или сокета в таблице процесса. Отдельное имя
типа делает интерфейс \texttt{os.h} понятным и позволяет при необходимости
изменить представление дескриптора, не меняя код родителя.

\subsection{Тип \texttt{ssize\_t}}

Знаковый целочисленный тип для размеров. Отличается от \texttt{size\_\allowbreak{}t} наличием
знака, что необходимо, так как функции \texttt{read}, \texttt{write},
\texttt{readlink} возвращают число байт либо \texttt{-1} при ошибке.

\subsection{Функция \texttt{pipe2}}

\begin{lstlisting}
int pipe2(int fds[2], int flags);
\end{lstlisting}

Создаёт неименованный канал~--- однонаправленный буфер в ядре. В \texttt{fds[0]}
возвращается дескриптор для чтения, в \texttt{fds[1]}~--- для записи. Данные,
записанные в \texttt{fds[1]}, можно прочитать из \texttt{fds[0]} в порядке FIFO.
В отличие от обычной \texttt{pipe}, функция \texttt{pipe2} принимает флаги и
позволяет атомарно установить \texttt{O\_\allowbreak{}CLOEXEC}. Это важно в многопоточных
программах, где между \texttt{pipe} и \texttt{fcntl} мог бы произойти
\texttt{fork}+\texttt{exec}.

\subsection{Функция \texttt{open} и её флаги}

\begin{lstlisting}
int fd = open(path, O_WRONLY | O_CREAT | O_TRUNC | O_CLOEXEC, 0644);
\end{lstlisting}

Флаги объединяются побитовым ИЛИ:

\begin{itemize}[noitemsep]
  \item \texttt{O\_\allowbreak{}WRONLY}~--- открыть файл только на запись;
  \item \texttt{O\_\allowbreak{}CREAT}~--- создать файл, если он не существует;
  \item \texttt{O\_\allowbreak{}TRUNC}~--- если файл существует, обнулить его содержимое;
  \item \texttt{O\_\allowbreak{}CLOEXEC}~--- закрыть дескриптор автоматически при вызове \texttt{exec}.
\end{itemize}

Число \texttt{0644} (восьмеричное)~--- права нового файла: владелец читает и
пишет, остальные только читают.

\subsection{Флаг \texttt{O\_CLOEXEC}}

Флаг <<close on exec>>. Без него дочерняя программа после \texttt{exec}
унаследовала бы \textit{все} открытые дескрипторы родителя, в том числе
записывающие концы каналов. Тогда \texttt{child} никогда не получил бы
конец файла (EOF) на \texttt{stdin}, потому что у него самого оставалась бы
открытая копия записывающего конца, и программа зависла бы. Благодаря
\texttt{O\_\allowbreak{}CLOEXEC} после \texttt{exec} остаются только дескрипторы 0, 1, 2,
явно установленные через \texttt{dup2}.

\subsection{Функция \texttt{fork}}

Создаёт копию текущего процесса. В родителе возвращает идентификатор
(pid) ребёнка, в ребёнке~--- \texttt{0}, при ошибке~--- \texttt{-1}.
Это позволяет в одном коде развести поведение родителя и ребёнка:

\begin{lstlisting}
pid_t pid = fork();
if (pid < 0)  { /* ошибка */ }
if (pid == 0) { /* код ребёнка */ }
proc->pid = pid; /* код родителя */
\end{lstlisting}

\subsection{Функция \texttt{dup2}}

\begin{lstlisting}
dup2(in,  STDIN_FILENO);   /* stdin  ребёнка = канал */
dup2(out, STDOUT_FILENO);  /* stdout ребёнка = файл  */
\end{lstlisting}

\texttt{dup2(old, new)} делает дескриптор \texttt{new} копией \texttt{old},
при необходимости предварительно закрывая \texttt{new}. Именно так выполняется
\textbf{перенаправление ввода-вывода}: дочерняя программа не знает о каналах
и файлах, она просто читает из дескриптора 0 и пишет в дескриптор 1.
Важное свойство: копия, созданная \texttt{dup2}, \textbf{не} наследует
\texttt{O\_\allowbreak{}CLOEXEC}, поэтому переживёт \texttt{exec}.

\subsection{Функция \texttt{execl}}

Заменяет образ текущего процесса новой программой. Если вызов удался,
управление в старый код никогда не возвращается; возврат означает ошибку.
После неудачного \texttt{execl} в ребёнке вызывается \texttt{\_\allowbreak{}exit(127)}
(именно \texttt{\_\allowbreak{}exit}, а не \texttt{exit}, чтобы не сбрасывать копии
буферов stdio, унаследованные от родителя).

\subsection{Функция \texttt{readlink} и путь к \texttt{child}}

Файл \texttt{/\allowbreak{}proc/\allowbreak{}self/\allowbreak{}exe}~--- символическая ссылка на исполняемый файл
текущего процесса. Функция \texttt{readlink} читает её значение, то есть путь
вида \texttt{/\allowbreak{}.../\allowbreak{}build/\allowbreak{}bin/\allowbreak{}parent}. Затем \texttt{strrchr(buf,'/\allowbreak{}')} находит
последний слэш, после него подставляется имя \texttt{child}. Так
\texttt{parent} находит \texttt{child} в своей папке независимо от
текущей рабочей директории.

\subsection{Функция \texttt{write} и цикл \texttt{os\_write\_all}}

\texttt{write} может записать меньше байт, чем запрошено (например, при
заполненном канале или прерывании сигналом). Поэтому запись выполняется в
цикле до тех пор, пока не будут переданы все данные. Код ошибки
\texttt{EINTR} (вызов прерван сигналом) не считается ошибкой~--- запись
просто повторяется.

\subsection{Функция \texttt{waitpid}}

\begin{lstlisting}
waitpid(proc->pid, &status, 0);
\end{lstlisting}

Блокирует родителя до завершения указанного дочернего процесса и
записывает в \texttt{status} причину завершения. Помимо получения кода
возврата, вызов \textbf{освобождает запись о процессе в ядре}; пока
родитель не вызвал \texttt{waitpid}, завершившийся ребёнок остаётся
<<зомби-процессом>>.

\subsection{Макросы \texttt{WIFEXITED}, \texttt{WEXITSTATUS}, \texttt{WIFSIGNALED}, \texttt{WTERMSIG}}

Значение \texttt{status} упаковано в битовое поле, поэтому расшифровывается
макросами:

\begin{center}
\begin{tabular}{@{}p{4.3cm}p{10cm}@{}}
\toprule
\textbf{Макрос} & \textbf{Смысл}\\
\midrule
\texttt{WIFEXITED(s)} & Истина, если процесс завершился нормально (\texttt{return} из \texttt{main} или \texttt{exit}).\\
\texttt{WEXITSTATUS(s)} & Код возврата (имеет смысл, только если \texttt{WIFEXITED} истинно).\\
\texttt{WIFSIGNALED(s)} & Истина, если процесс был завершён сигналом (например, \texttt{SIGSEGV}, \texttt{SIGKILL}).\\
\texttt{WTERMSIG(s)} & Номер сигнала, убившего процесс.\\
\bottomrule
\end{tabular}
\end{center}

\subsection{Сигнал \texttt{SIGPIPE}}

Если записать в канал, у которого не осталось читателей, ядро посылает
процессу \texttt{SIGPIPE}, что по умолчанию убивает программу. В родителе
вызывается \texttt{signal(SIGPIPE, SIG\_\allowbreak{}IGN)}: тогда \texttt{write} просто
вернёт ошибку \texttt{EPIPE}, которую можно обработать и корректно
сообщить пользователю.

\section{Платформонезависимые функции}

\subsection{\texttt{read\_line}}

Читает строку произвольной длины из \texttt{FILE*}. Буфер начинается с 128
байт и удваивается через \texttt{realloc}, когда заканчивается место. Это
исключает типичную ошибку фиксированного буфера (\texttt{char buf[256]}),
из-за которой длинная строка обрезалась бы или портила память. Параметр
\texttt{keep\_\allowbreak{}nl} определяет, сохранять ли \texttt{\textbackslash n}; символ
\texttt{\textbackslash r} перед \texttt{\textbackslash n} (формат CRLF)
отбрасывается. Ошибка передаётся через \texttt{*err}, а различие между
<<конец ввода>> и <<ошибка>> делается проверкой \texttt{ferror}.

\subsection{Инверсия строки в \texttt{child}}

\begin{lstlisting}
static void reverse(char *s, size_t len)
{
    size_t i, j;
    char tmp;
    if (len < 2) return;
    for (i = 0, j = len - 1; i < j; i++, j--) {
        tmp = s[i]; s[i] = s[j]; s[j] = tmp;
    }
}
\end{lstlisting}

Алгоритм «двух указателей»: один идёт с начала, другой с конца, символы
меняются местами до встречи в середине. Проверка \texttt{len < 2} необходима,
так как \texttt{len} имеет беззнаковый тип \texttt{size\_\allowbreak{}t}, и выражение
\texttt{len - 1} для пустой строки дало бы огромное число. После каждой
строки в \texttt{child} вызывается \texttt{fflush(stdout)}, чтобы данные
попадали в файл сразу, а не копились в буфере.

\subsection{Обработка ошибок}

\texttt{err\_\allowbreak{}errno} печатает сообщение вида
\texttt{<что делали>: <strerror(errno)>}. \texttt{err\_\allowbreak{}sys} в Linux
вызывает \texttt{err\_\allowbreak{}errno}. Все функции слоя \texttt{os} возвращают
\texttt{0} или \texttt{-1}, а сообщение об ошибке выводится на месте возникновения.

\section{Алгоритм работы родительского процесса}

\begin{enumerate}
  \item Определить путь к \texttt{child} (\texttt{os\_\allowbreak{}sibling\_\allowbreak{}path}).
  \item Прочитать два имени файлов.
  \item Открыть файлы на запись и создать два канала.
  \item Запустить \texttt{child1} (\texttt{stdin} = \texttt{pipe1}, \texttt{stdout} = файл~1).
        Закрыть в родителе читающий конец канала и дескриптор файла~1.
  \item Запустить \texttt{child2} аналогично, закрыть его читающий конец и файл~2.
  \item В цикле читать строки из \texttt{stdin}: нечётные записывать в
        \texttt{pipe1}, чётные в \texttt{pipe2}; при отсутствии завершающего
        \texttt{\textbackslash n} добавлять его.
  \item По концу ввода (Ctrl+D) закрыть
        записывающие концы каналов~--- дети получат EOF.
  \item Дождаться завершения обоих детей и проверить их коды возврата.
\end{enumerate}

Весь код очистки реализован через метку \texttt{cleanup} и \texttt{goto}:
независимо от места возникновения ошибки выполняется один и тот же
набор закрытий и ожиданий, что предотвращает утечки дескрипторов.
Отдельные <<закрыть в родителе>> после запуска ребёнка нужны, чтобы
в родителе не осталось лишних копий дескрипторов (иначе канал
не получит EOF).


\begin{itemize}
  \item \textbf{Ошибка \texttt{CMakeCache.txt directory is different}.}
        Папка \texttt{build} была создана CMake в другом окружении (другие
        пути и генератор), а затем использовалась из WSL
        (путь \texttt{/\allowbreak{}mnt/\allowbreak{}c/\allowbreak{}...}). Кэш конфигурации привязан к абсолютным
        путям и генератору, поэтому после смены окружения папку
        \texttt{build} нужно удалить и создать заново.
  \item \textbf{Ошибка \texttt{CMAKE\_\allowbreak{}C\_\allowbreak{}COMPILER not set}.} В WSL не были
        установлены компилятор и \texttt{make}. Решение:
        \texttt{sudo apt install build-essential cmake}.
\end{itemize}

\section{Выводы}

В ходе работы были изучены и применены на практике:

\begin{enumerate}
  \item \textbf{Модель процессов.} Освоена пара \texttt{fork}/\texttt{exec}
        в Linux; понято, что создание процесса (\texttt{fork}) и запуск в нём
        новой программы (\texttt{exec}) являются двумя отдельными шагами.
  \item \textbf{Неименованные каналы.} Канал~--- односторонний буфер в
        ядре. Однонаправленность и семантика EOF (конец данных наступает
        только когда закрыты \textit{все} записывающие концы) определяют
        правильное закрытие дескрипторов в родителе и ребёнке.
  \item \textbf{Перенаправление потоков.} С помощью \texttt{dup2}
        стандартные потоки ребёнка
        подменяются каналом и файлом, поэтому \texttt{child} является обычной
        консольной программой и не зависит от способа запуска.
  \item \textbf{Наследование дескрипторов.} Утечка дескрипторов в
        дочерний процесс~--- частая причина зависаний. Для её
        предотвращения применяется флаг \texttt{O\_\allowbreak{}CLOEXEC}.
  \item \textbf{Ожидание процессов.} Вызов \texttt{waitpid} необходим, чтобы получить результат
        работы ребёнка и не оставлять зомби-процессов. Разбор статуса
        через \texttt{WIFEXITED}/\texttt{WIFSIGNALED} позволяет отличить
        нормальное завершение от аварийного.
  \item \textbf{Надёжный ввод-вывод.} \texttt{write} может записать не
        все данные, а системные вызовы могут быть прерваны сигналом
        (\texttt{EINTR}), поэтому используется цикл записи и повтор вызова.
        Чтение строк произвольной длины реализовано динамическим буфером.
  \item \textbf{Модульность.} Разделение кода на общий слой, интерфейс
        \texttt{os.h} и его реализацию \texttt{os\_\allowbreak{}linux.c} позволило отделить
        логику родителя от системных вызовов и упростить чтение и отладку
        программы.
  \item \textbf{Инструменты.} Освоена система сборки CMake (статические
        библиотеки, зависимости целей, \texttt{ctest}), а также
        особенности запуска проекта в WSL и проблема кэша CMake при смене
        среды сборки.
\end{enumerate}

Практически важный вывод: большинство ошибок при работе с процессами и
каналами связано не с самими вызовами, а с \textbf{владением
дескрипторами}~--- кто и когда закрывает какой конец канала. Именно
аккуратное закрытие лишних копий обеспечивает корректное завершение
программы.

\medskip
\noindent\textbf{Личные результаты.} В ходе выполнения лабораторной работы я
научился создавать процессы с помощью системного вызова \texttt{fork} и
запускать в них другие программы через \texttt{exec}. Я разобрался, как
родительский процесс управляет дочерними: передаёт им данные через
неименованные каналы (\texttt{pipe}), перенаправляет их ввод и вывод
с помощью \texttt{dup2}, ожидает завершения через \texttt{waitpid} и по
статусу определяет, завершился процесс нормально или был убит сигналом.
Также я понял, как важно закрывать ненужные дескрипторы и как работают
наследование дескрипторов, конец файла (EOF) в канале и сигналы
(например, \texttt{SIGPIPE} и \texttt{SIGSEGV}). Кроме того, я
научился находить и исправлять ошибки при работе с процессами, в том числе
падение дочернего процесса из-за неподключённых заголовочных файлов.


\end{document}